\section*{第 \ref{ch:g6:rocks-and-minerals} 章——岩石与矿物}

\begin{solution}{exo:g6:rocks-and-minerals:1}
花岗岩：\omterm{def:g6:rocks-and-minerals:rock-families}{火成岩}。砂岩：\omterm{def:g6:rocks-and-minerals:rock-families}{沉积岩}。大理岩：\omterm{def:g6:rocks-and-minerals:rock-families}{变质岩}。玄武岩：
\omterm{def:g6:rocks-and-minerals:rock-families}{火成岩}。石灰岩：\omterm{def:g6:rocks-and-minerals:rock-families}{沉积岩}。板岩：\omterm{def:g6:rocks-and-minerals:rock-families}{变质岩}。
\end{solution}

\begin{solution}{exo:g6:rocks-and-minerals:2}
石英是\omterm{def:g6:rocks-and-minerals:mineral}{矿物}，花岗岩是\omterm{def:g3:stones-pebbles-sand:rock}{岩石}。\omterm{def:g6:rocks-and-minerals:mineral}{矿物}是由原子按规则方式排列而成的单一天然
物质；\omterm{def:g3:stones-pebbles-sand:rock}{岩石}则由一种或几种\omterm{def:g6:rocks-and-minerals:mineral}{矿物}紧密结合而成
（花岗岩含长石、石英和云母）。
\end{solution}

\begin{solution}{exo:g6:rocks-and-minerals:3}
滑石（1）、方解石（3）、石英（7）、金刚石（10）。金刚石最硬，
能划伤其他所有\omterm{def:g6:rocks-and-minerals:mineral}{矿物}。
\end{solution}

\begin{solution}{exo:g6:rocks-and-minerals:4}
“熔融\omterm{def:g3:stones-pebbles-sand:rock}{岩石}冷却后成为\omterm{def:g6:rocks-and-minerals:rock-families}{火成岩}。石灰岩在
热和压力作用下成为大理岩。花岗岩经
\omterm{def:g5:water-shapes-the-land:weathering}{风化}成为\omterm{def:g3:stones-pebbles-sand:sand}{沙粒}。”
\end{solution}

\begin{solution}{exo:g6:rocks-and-minerals:5}
(a) 玄武岩。(b) 大理岩。(c) 砂岩。
\end{solution}

\begin{solution}{exo:g6:rocks-and-minerals:6}
其\omterm{def:g6:rocks-and-minerals:hardness}{硬度}大于指甲（约 2.5），小于
铜币（约 3.5）。它可能是方解石（3）。
\end{solution}

\begin{solution}{exo:g6:rocks-and-minerals:7}
石英：$21 \div 60 = 0.35$，即 35~\%。长石：$33 \div 60 = 0.55$，
即 55~\%。云母：$6 \div 60 = 0.10$，即 10~\%。检验：
$35 + 55 + 10 = 100$。
\end{solution}

\begin{solution}{exo:g6:rocks-and-minerals:8}
花岗岩在地下深处冷却非常缓慢，其晶粒有充足
时间生长而变大。玄武岩来自\omterm{def:g4:volcanoes-earthquakes-first-look:eruption}{熔岩}，\omterm{def:g4:volcanoes-earthquakes-first-look:eruption}{熔岩}在地表
迅速冷却，晶粒没有生长时间，小得肉眼难以分辨。
\end{solution}

\begin{solution}{exo:g6:rocks-and-minerals:9}
(a) 石灰岩是\omterm{def:g6:rocks-and-minerals:rock-families}{沉积岩}；大理岩是\omterm{def:g6:rocks-and-minerals:rock-families}{变质岩}。(b) 在地下深处，
热和压力使其变化而未熔化：其方解石生长成相互嵌合的
闪亮晶体。(c) 两者都由同一种\omterm{def:g6:rocks-and-minerals:mineral}{矿物}方解石组成，
遇酸冒泡的正是方解石。
\end{solution}

\begin{solution}{exo:g6:rocks-and-minerals:10}
(a) 能：石英（7）比玻璃（约 5.5）硬。(b) 不能：萤石（4）
比玻璃软。(c) 能：钢钉（约 6.5）比
长石（6）硬。
\end{solution}

\begin{solution}{exo:g6:rocks-and-minerals:11}
例如：(1) 山上的花岗岩（\omterm{def:g6:rocks-and-minerals:rock-families}{火成岩}）经\omterm{def:g5:water-shapes-the-land:weathering}{风化}破碎，
其石英颗粒被释放出来；(2) \omterm{def:g1:landscapes-around-us:water}{河流}将其带走：\omterm{def:g5:water-shapes-the-land:erosion}{侵蚀}；
(3) 颗粒在海底一层层\omterm{def:g5:water-shapes-the-land:deposition}{沉积}：\omterm{def:g5:water-shapes-the-land:deposition}{沉积}；
(4) 颗粒被其他层覆盖后受压并胶结成砂岩：固结
（此时\omterm{def:g3:stones-pebbles-sand:rock}{岩石}为\omterm{def:g6:rocks-and-minerals:rock-families}{沉积岩}）；
(5) 砂岩被带到地下深处，在热和压力作用下
变成\omterm{def:g6:rocks-and-minerals:rock-families}{变质岩}。
\end{solution}

\begin{solution}{exo:g6:rocks-and-minerals:12}
(a) 钢（约 6.5）划不动它，且它能划伤玻璃：\omterm{def:g6:rocks-and-minerals:hardness}{硬度}
大于 6.5。黄玉（8）能划伤它：\omterm{def:g6:rocks-and-minerals:hardness}{硬度}小于 8。其\omterm{def:g6:rocks-and-minerals:hardness}{硬度}为 7，即石英。
(b) 石英比与之摩擦的几乎所有物质（其他颗粒、\omterm{def:g3:stones-pebbles-sand:pebble}{卵石}）都硬，
几乎不会被磨损。方解石（3）和云母
（2 至 3）则软得多，摩擦会将它们磨成细小碎屑
并被带走。
\end{solution}

\begin{solution}{pb:g6:rocks-and-minerals:1}
\textbf{1.} 花岗岩，\omterm{def:g6:rocks-and-minerals:rock-families}{火成岩}；其\omterm{def:g6:rocks-and-minerals:mineral}{矿物}为长石（粉红色或
白色）、石英（灰色，玻璃光泽）和云母（黑色，有光泽）。

\textbf{2.} 石灰岩，\omterm{def:g6:rocks-and-minerals:rock-families}{沉积岩}。贝壳属于生活在海中的
动物；它们的遗骸\omterm{def:g5:water-shapes-the-land:deposition}{沉积}下来并固结成
\omterm{def:g3:stones-pebbles-sand:rock}{岩石}。

\textbf{3.} 板岩，\omterm{def:g6:rocks-and-minerals:rock-families}{变质岩}，由页岩变质而成。

\textbf{4.} 大理岩，\omterm{def:g6:rocks-and-minerals:rock-families}{变质岩}，由石灰岩变质而成。

\textbf{5.} 石英：钢（约 6.5）划不动它，它能划伤
玻璃（约 5.5）；其\omterm{def:g6:rocks-and-minerals:hardness}{硬度}为 7。

\textbf{6.} 能被钢（约 6.5）划伤，但不能被玻璃
（约 5.5）划伤：\omterm{def:g6:rocks-and-minerals:hardness}{硬度}介于 5.5 与 6.5 之间。它是长石（正长石，6）。

\textbf{7.} 云母（\omterm{def:g6:rocks-and-minerals:hardness}{硬度} 2 至 3）。

\textbf{8.} 能。铺路的是石灰岩，由方解石（3）组成，比
铜币（约 3.5）软。

\textbf{9.} \omterm{def:g5:water-shapes-the-land:weathering}{风化}：雨水缓慢磨损石头表面，
使石面逐渐被磨去。

\textbf{10.} 山体经\omterm{def:g5:water-shapes-the-land:weathering}{风化}，其石英颗粒被\omterm{def:g1:landscapes-around-us:water}{河流}
带走，并在湖底或海底一层层\omterm{def:g5:water-shapes-the-land:deposition}{沉积}；经掩埋
和固结后，可形成砂岩，即一种\omterm{def:g6:rocks-and-minerals:rock-families}{沉积岩}。

\textbf{11.} 三类\omterm{def:g3:stones-pebbles-sand:rock}{岩石}都有：\omterm{def:g6:rocks-and-minerals:rock-families}{火成岩}（喷泉中的花岗岩）、
\omterm{def:g6:rocks-and-minerals:rock-families}{沉积岩}（铺路的石灰岩）、\omterm{def:g6:rocks-and-minerals:rock-families}{变质岩}（屋顶的板岩
和雕像的大理岩）。

\textbf{12.} 最硬的\omterm{def:g6:rocks-and-minerals:mineral}{矿物}是花岗岩中的石英：莫氏
\omterm{def:g6:rocks-and-minerals:hardness}{硬度}为 \textbf{7}。
\end{solution}
